\begin{onehalfspace}

\section{Qiskit: Il Framework IBM per il Calcolo Quantistico}

\textbf{Qiskit} è un framework open-source sviluppato da \ac{IBM} per la programmazione di computer quantistici \cite{qiskit2023}.\footnote{Qiskit (da \textit{Quantum Information Software Kit}) fu rilasciato pubblicamente da IBM nel marzo 2017 come progetto open-source su GitHub. È diventato il framework quantistico più diffuso al mondo, con oltre 550.000 utenti registrati sulla piattaforma IBM Quantum Experience.} Scritto in Python, consente di costruire, simulare ed eseguire circuiti quantistici su hardware reale attraverso IBM Quantum Experience. L'architettura di Qiskit si articola in quattro livelli:
\begin{itemize}
    \item \textbf{Terra}: costruzione e ottimizzazione di circuiti quantistici;
    \item \textbf{Aer}: simulatore ad alte prestazioni (statevector, QASM, noise model);
    \item \textbf{Ignis}: caratterizzazione dell'hardware e mitigazione del rumore;
    \item \textbf{Aqua}: algoritmi quantistici ad alto livello.
\end{itemize}

\section{Implementazione dell'Algoritmo di Shor con Qiskit}

Le implementazioni sviluppate in questo capitolo traggono ispirazione dal repository open-source di riferimento \cite{github_shor}, adattato e esteso per includere l'analisi del rumore e i moduli di ottimizzazione basati su intelligenza artificiale.

\subsection{Quantum Fourier Transform}

La \acf{QFT} è la subroutine centrale dell'algoritmo di Shor. Trasforma gli stati della base computazionale come:
\begin{equation}
    \text{QFT}\ket{j} = \frac{1}{\sqrt{N}} \sum_{k=0}^{N-1} e^{2\pi i jk/N} \ket{k}
\end{equation}
Il circuito richiede $O(n^2)$ gate, dove $n$ è il numero di qubit. L'implementazione in Qiskit sfrutta la struttura ricorsiva della QFT tramite porte di Hadamard e rotazioni di fase controllate:

\begin{lstlisting}[caption={Implementazione della Quantum Fourier Transform}]
from qiskit import QuantumCircuit
import numpy as np

def qft_circuit(n_qubits):
    """Quantum Fourier Transform su n qubit."""
    qc = QuantumCircuit(n_qubits, name='QFT')
    for j in range(n_qubits):
        qc.h(j)
        for k in range(j + 1, n_qubits):
            angle = np.pi / (2 ** (k - j))
            qc.cp(angle, k, j)
    # SWAP finale per correggere l'ordine dei bit
    for i in range(n_qubits // 2):
        qc.swap(i, n_qubits - i - 1)
    return qc

def inverse_qft(n_qubits):
    """QFT inversa: coniugato trasposto della QFT."""
    qc = qft_circuit(n_qubits)
    return qc.inverse()
\end{lstlisting}

\subsection{Modular Exponentiation}

Il blocco di \textit{modular exponentiation} implementa l'operazione $U_f\ket{x}\ket{y} = \ket{x}\ket{y \cdot a^x \bmod N}$. Per istanze concrete, questo blocco viene compilato in gate elementari. Di seguito, l'implementazione per il caso $N = 15$, $a = 7$:

\begin{lstlisting}[caption={Modular exponentiation per N=15, a=7}]
from qiskit import QuantumCircuit

def mod_exp_15_7(power):
    """
    Implementa U^power dove U|y> = |7y mod 15>.
    Valida per power in {1, 2, 4, 8}.
    """
    U = QuantumCircuit(4, name=f'U^{power}')
    # Permutazioni dei computational basis states per 7^power mod 15
    # 7^1 mod 15 = 7:  1->7, 7->4, 4->13, 13->1
    # 7^2 mod 15 = 4:  1->4, 4->1, 7->13, 13->7
    if power % 4 == 1:
        U.swap(0, 1)
        U.swap(1, 2)
        U.swap(2, 3)
        U.x(0); U.x(2)
    elif power % 4 == 2:
        U.swap(1, 3)
        U.swap(0, 2)
    elif power % 4 == 0:
        pass  # Identita per U^4 = I (mod 15 con a=7)
    return U
\end{lstlisting}

Gli operatori di moltiplicazione modulare si costruiscono efficientemente con porte SWAP, sfruttando la struttura di permutazione dell'operazione $x \mapsto ax \bmod N$. La Figura~\ref{fig:mod_circuits} mostra i circuiti per $M_2$ e $M_4$ modulo 15.

\begin{figure}[H]
\centering
\subfloat[Operatore $M_2 \pmod{15}$]{\includegraphics[width=0.48\textwidth]{01_M2_circuit.pdf}}
\hfill
\subfloat[Operatore $M_4 \pmod{15}$]{\includegraphics[width=0.48\textwidth]{04_M4_circuit.pdf}}
\caption{Circuiti degli operatori di moltiplicazione modulare per $N=15$, costruiti con porte SWAP. $M_2$ richiede 3 SWAP in cascata; $M_4$ richiede 2 SWAP paralleli. Entrambi vengono decomposti in gate CNOT per l'esecuzione su hardware reale.}
\label{fig:mod_circuits}
\end{figure}

\subsection{Circuito Completo dell'Algoritmo di Shor}

Il circuito completo dell'algoritmo di Shor per $N = 15$, $a = 7$, utilizza due registri: un registro di \textit{counting} (per la phase estimation) e un registro di \textit{lavoro} inizializzato a $\ket{1}$:

\begin{lstlisting}[caption={Circuito completo di Shor per N=15}]
from qiskit import QuantumCircuit, transpile
from qiskit_aer import AerSimulator
from math import gcd

def shor_circuit_N15(a=7, n_count=8):
    """
    Algoritmo di Shor per N=15.
    n_count: numero di qubit nel registro di counting.
    """
    n_work = 4  # qubit di lavoro (log2(15) + 1)
    qc = QuantumCircuit(n_count + n_work, n_count)

    # Passo 1: Hadamard sul registro di counting
    for q in range(n_count):
        qc.h(q)

    # Passo 2: Inizializza il registro di lavoro a |1>
    qc.x(n_count)

    # Passo 3: Controlled-U^(2^j) per j = 0, ..., n_count-1
    for j in range(n_count):
        power = 2 ** j
        ctrl_U = mod_exp_15_7(power).control(1)
        qc.append(ctrl_U, [j] + list(range(n_count, n_count + n_work)))

    qc.barrier()

    # Passo 4: QFT inversa sul registro di counting
    qc.append(inverse_qft(n_count), range(n_count))

    # Passo 5: Misura del registro di counting
    qc.measure(range(n_count), range(n_count))
    return qc

# Esecuzione su simulatore ideale
qc = shor_circuit_N15(a=7, n_count=8)
sim = AerSimulator()
job = sim.run(transpile(qc, sim), shots=1024)
counts = job.result().get_counts()
print(counts)
# Output atteso: picchi a multipli di 2^n_count / r
# con r=4 (periodo di 7 mod 15): picchi a 0, 64, 128, 192
\end{lstlisting}

La struttura del circuito è visibile in Figura~\ref{fig:qpe_circuit_impl}: il registro di controllo viene posto in sovrapposizione, entanglato tramite le controlled-$U$ di moltiplicazione modulare, e infine misurato dopo la QFT$^{-1}$.

\begin{figure}[H]
\centering
\includegraphics[width=0.95\textwidth]{09_QPE_circuit_N15.pdf}
\caption{Circuito QPE completo per $N=15$, $a=2$, generato con Qiskit (8 qubit di controllo, 4 qubit di lavoro). I blocchi $M_2$ e $M_4$ rappresentano gli operatori di moltiplicazione modulare controllata; la QFT$^{-1}$ finale estrae l'informazione di fase dal registro di controllo.}
\label{fig:qpe_circuit_impl}
\end{figure}

\subsection{Post-Processing: Estrazione del Fattore}

Dopo la misura del registro di counting, il valore $y$ ottenuto viene usato per stimare la frazione $y / 2^{n\_count} \approx s/r$ tramite l'algoritmo delle frazioni continue, ottenendo il periodo $r$. Se $r$ è pari e $a^{r/2} \not\equiv -1 \pmod{N}$, allora $\gcd(a^{r/2} \pm 1, N)$ fornisce un fattore non banale di $N$:

\begin{lstlisting}[caption={Post-processing: dalle misure ai fattori}]
from fractions import Fraction

def extract_period(measured_value, n_count, N):
    """Stima il periodo r dalla misura tramite frazioni continue."""
    phase = measured_value / (2 ** n_count)
    frac = Fraction(phase).limit_denominator(N)
    r = frac.denominator
    return r

def find_factors(a, r, N):
    """Calcola i fattori di N da a e il periodo r."""
    if r % 2 != 0:
        return None, None
    candidate1 = gcd(a ** (r // 2) - 1, N)
    candidate2 = gcd(a ** (r // 2) + 1, N)
    if candidate1 not in [1, N]:
        return candidate1, N // candidate1
    if candidate2 not in [1, N]:
        return candidate2, N // candidate2
    return None, None

# Esempio: measured_value = 64 con n_count = 8, N = 15, a = 7
r = extract_period(64, 8, 15)    # r = 4
p, q = find_factors(7, 4, 15)   # p=3, q=5
print(f"Fattori di 15: {p} x {q}")   # 3 x 5
\end{lstlisting}

\section{Analisi Sperimentale su Simulatore Rumoroso}

\subsection{Modello di Rumore}

Per simulare le condizioni di un processore NISQ reale, si utilizza il simulatore Aer di Qiskit con un modello di rumore che include:
\begin{itemize}
    \item Rumore depolarizzante su gate a singolo qubit ($\epsilon_1 = 10^{-3}$);
    \item Rumore depolarizzante su gate a due qubit ($\epsilon_2 = 10^{-2}$);
    \item Amplitude damping con $T_1 = 100\ \mu\text{s}$ e dephasing con $T_2 = 80\ \mu\text{s}$.
\end{itemize}

\begin{lstlisting}[caption={Configurazione del modello di rumore per Shor}]
from qiskit_aer.noise import NoiseModel, depolarizing_error, thermal_relaxation_error

noise_model = NoiseModel()

# Errori su gate a singolo qubit
error_1q = depolarizing_error(1e-3, 1)
noise_model.add_all_qubit_quantum_error(error_1q, ['h', 'x', 'rz', 'cp'])

# Errori su gate a due qubit
error_2q = depolarizing_error(1e-2, 2)
noise_model.add_all_qubit_quantum_error(error_2q, ['cx', 'swap'])

# Rilassamento termico (T1, T2)
t1, t2 = 100e3, 80e3  # in ns
gate_time = 50        # ns per gate a singolo qubit
error_relax = thermal_relaxation_error(t1, t2, gate_time)
noise_model.add_all_qubit_quantum_error(error_relax, ['h', 'x'])
\end{lstlisting}

\subsection{Tecniche di Soppressione del Rumore: Dynamical Decoupling e Gate Twirling}

Oltre al modello di rumore per la simulazione, il tutorial ufficiale IBM \cite{ibm_shor_tutorial} adotta due tecniche di soppressione del rumore che possono essere applicate anche in fase di esecuzione su hardware reale o su simulatori che le supportano:

\begin{itemize}
    \item \textbf{Dynamical Decoupling (sequenza XpXm)}: sequenze di pulse $X$-$X$ applicate ai qubit inattivi durante i gate lunghi, in modo da mediare a zero l'interazione con l'ambiente. Su \texttt{ibm\_marrakesh}, questa tecnica viene abilitata tramite le opzioni di transpilazione di Qiskit Runtime;
    \item \textbf{Gate Twirling}: randomizzazione dei gate a due qubit tramite frame casuali di Pauli, che converte errori coerenti in errori stocastici, rendendo il rumore più prevedibile e la mitigazione più efficace. Su hardware IBM, il Gate Twirling riduce la varianza degli errori di gate prima dell'applicazione di ZNE o PEC.
\end{itemize}

A titolo di riferimento, il circuito di Shor per $N=15$, $a=2$, compilato su \texttt{ibm\_marrakesh} con optimization level 2, produce una profondità di \textbf{187 gate a due qubit} e un conteggio totale di \textbf{260 gate a due qubit} -- valori che illustrano concretamente il costo circuitale di un'istanza piccola ma realistica \cite{ibm_shor_tutorial}.

\subsection{Risultati su Simulatore Ideale vs.\ Rumoroso}

Eseguendo il circuito di Shor per $N = 15$, $a = 7$ con $n_{\text{count}} = 8$ qubit di counting si ottengono i seguenti risultati:

\begin{itemize}
    \item \textbf{Simulatore ideale}: picchi netti a $y \in \{0, 64, 128, 192\}$ (multipli di $256/4$), con probabilità totale di misurare un valore utile $> 95\%$.
    \item \textbf{Simulatore rumoroso} (senza mitigazione): la distribuzione si appiattisce, con le misure utili che raggiungono una probabilità del $40$--$60\%$.
    \item \textbf{Con ZNE}: il recupero delle misure utili sale al $65$--$75\%$.
    \item \textbf{Con ZNE + readout mitigation}: il recupero sale al $75$--$85\%$.
\end{itemize}

La Figura~\ref{fig:histograms_N15} mostra la distribuzione delle misure del registro di controllo nei due scenari: simulatore ideale e hardware reale \texttt{ibm\_marrakesh} (con dynamical decoupling e gate twirling abilitati).

\begin{figure}[H]
\centering
\subfloat[Simulatore ideale (1024 shot)]{\includegraphics[width=0.48\textwidth]{10_histogram_simulator_N15.pdf}}
\hfill
\subfloat[Hardware reale \texttt{ibm\_marrakesh} (1024 shot)]{\includegraphics[width=0.48\textwidth]{11_histogram_hardware_N15.pdf}}
\caption{Distribuzione delle misure del registro di controllo per $N=15$, $a=2$. I quattro picchi a $\{0, 64, 128, 192\}$ corrispondono alle fasi $k/r$ con $r=4$ ($k \in \{0,1,2,3\}$). Sul simulatore ideale i picchi sono netti e bilanciati; sull'hardware reale emerge un fondo di rumore che viene filtrato statisticamente con una soglia al 50\% del picco massimo.}
\label{fig:histograms_N15}
\end{figure}

\section{Ottimizzazione mediante Modelli di Intelligenza Artificiale}

\subsection{Motivazione}

Le tecniche di mitigazione tradizionali (ZNE, PEC, readout mitigation) trattano il rumore in modo generico. I modelli di intelligenza artificiale offrono una prospettiva complementare: apprendendo il comportamento specifico dell'hardware e del circuito, possono caratterizzare il rumore con maggiore precisione, ottimizzare i parametri circuitali e migliorare il post-processing dei risultati \cite{ml_mitigation2023, shor_hpec2023}.

\subsection{Caratterizzazione del Rumore con Machine Learning}

Un modello di regressione (es. Random Forest o rete neurale) può essere addestrato a predire la fedeltà del circuito di Shor in funzione dei parametri di rumore dell'hardware ($\epsilon_1$, $\epsilon_2$, $T_1$, $T_2$, profondità del circuito):

\begin{lstlisting}[caption={Regressore per la fedeltà del circuito di Shor}]
from sklearn.ensemble import RandomForestRegressor
import numpy as np

# Feature: [eps_1q, eps_2q, T1, T2, circuit_depth]
# Target: fidelity (probabilita' di successo della fattorizzazione)

# Dataset generato da simulazioni con parametri variabili
X_train = np.array([
    [1e-3, 1e-2, 100e3, 80e3, 350],
    [5e-4, 5e-3, 200e3, 150e3, 350],
    # ...
])
y_train = np.array([0.52, 0.74, ...])

model = RandomForestRegressor(n_estimators=100, random_state=42)
model.fit(X_train, y_train)

# Predire la fedelta' per un nuovo hardware
params = [[2e-3, 1.5e-2, 80e3, 60e3, 350]]
predicted_fidelity = model.predict(params)
print(f"Fedeltà predetta: {predicted_fidelity[0]:.3f}")
\end{lstlisting}

\subsection{Ottimizzazione dei Parametri Circuitali con Reinforcement Learning}

Un agente di Reinforcement Learning può apprendere una politica di ottimizzazione del circuito di Shor -- scegliendo, ad esempio, la decomposizione ottimale della modular exponentiation, il routing dei qubit e la scelta delle coppie di qubit con la migliore fedeltà -- massimizzando la probabilità di successo della fattorizzazione:

\begin{lstlisting}[caption={Struttura dell'ambiente RL per ottimizzazione circuitale}]
import gymnasium as gym
import numpy as np

class ShorOptimizationEnv(gym.Env):
    """
    Ambiente RL per l'ottimizzazione del circuito di Shor.
    Stato: parametri del circuito corrente + metriche di rumore.
    Azione: modifica di un parametro (routing, decomposizione).
    Reward: fedelta' misurata su simulatore rumoroso.
    """
    def __init__(self, N, a, noise_params):
        super().__init__()
        self.N = N
        self.a = a
        self.noise_params = noise_params
        self.action_space = gym.spaces.Discrete(10)
        self.observation_space = gym.spaces.Box(
            low=0, high=1, shape=(20,), dtype=np.float32
        )

    def step(self, action):
        # Applica azione (modifica parametro circuito)
        self._apply_action(action)
        # Esegui circuito su simulatore rumoroso
        fidelity = self._run_circuit()
        reward = fidelity
        terminated = fidelity > 0.90
        return self._get_obs(), reward, terminated, False, {}

    def _run_circuit(self):
        # Esegue il circuito di Shor e restituisce la fedelta'
        qc = shor_circuit_N15(self.a)
        sim = AerSimulator(noise_model=self._build_noise())
        job = sim.run(transpile(qc, sim), shots=512)
        counts = job.result().get_counts()
        useful = sum(v for k, v in counts.items()
                     if int(k, 2) % 64 == 0 and int(k, 2) > 0)
        return useful / 512
\end{lstlisting}

\subsection{Mitigazione Basata su Reti Neurali}

Le reti neurali possono essere addestrate a correggere le distribuzioni di probabilità prodotte dal circuito rumoroso, mappandole verso la distribuzione ideale. Questo approccio, noto come \textit{neural network-based error mitigation}, apprende la struttura dell'errore specifico del circuito di Shor e del processore utilizzato, superando i limiti delle tecniche di mitigazione standard che assumono modelli di rumore semplificati:

\begin{lstlisting}[caption={Rete neurale per la correzione delle distribuzioni di probabilità}]
import torch
import torch.nn as nn

class ErrorMitigationNet(nn.Module):
    """
    Mappa la distribuzione rumorosa p_noisy verso
    la distribuzione ideale p_ideal per il circuito di Shor.
    """
    def __init__(self, n_states):
        super().__init__()
        self.net = nn.Sequential(
            nn.Linear(n_states, 256),
            nn.ReLU(),
            nn.Linear(256, 256),
            nn.ReLU(),
            nn.Linear(256, n_states),
            nn.Softmax(dim=-1)
        )

    def forward(self, p_noisy):
        return self.net(p_noisy)

# Addestramento: coppie (p_noisy, p_ideal) generate su simulatore
n_states = 2 ** 8  # 256 stati per n_count = 8
model = ErrorMitigationNet(n_states)
optimizer = torch.optim.Adam(model.parameters(), lr=1e-3)
loss_fn = nn.KLDivLoss(reduction='batchmean')
\end{lstlisting}

\section{Risultati dell'Ottimizzazione AI-Assistita}

L'integrazione dei modelli di intelligenza artificiale con le tecniche di mitigazione tradizionali produce un miglioramento significativo della probabilità di successo della fattorizzazione rispetto al solo approccio classico:

\begin{itemize}
    \item \textbf{Senza mitigazione}: probabilità di successo $\approx 40$--$60\%$ su simulatore rumoroso.
    \item \textbf{ZNE + readout mitigation}: $\approx 75$--$85\%$.
    \item \textbf{RL-ottimizzato + ZNE}: $\approx 80$--$88\%$, grazie alla scelta ottimale del routing e della decomposizione.
    \item \textbf{Neural network mitigation}: $\approx 85$--$92\%$ su circuiti con parametri di rumore vicini al training set.
\end{itemize}

Questi risultati indicano che l'intelligenza artificiale può fornire un contributo concreto all'esecuzione dell'algoritmo di Shor su hardware NISQ, riducendo il gap tra simulatore ideale e hardware reale in modo complementare alle tecniche di mitigazione classiche.

\end{onehalfspace}
